% ============================================================
%  परिशिष्ट घ: परियोजना मूल्यांकन रूब्रिक (Project Assessment Rubrics)
% ============================================================
\chapter*{परिशिष्ट घ: परियोजना मूल्यांकन रूब्रिक}
\markboth{परियोजना मूल्यांकन रूब्रिक}{परियोजना मूल्यांकन रूब्रिक}
\addcontentsline{toc}{chapter}{परिशिष्ट घ: परियोजना मूल्यांकन रूब्रिक}

यह रूब्रिक शिक्षकों के लिए है ताकि वे पुस्तक की \textbf{परियोजनाओं (projects)} का
समान एवं निष्पक्ष मूल्यांकन कर सकें। प्रत्येक मानदंड (criterion) के लिए विद्यार्थी को
1--4 अंक दें; कुल 20 अंक।

\vspace{8pt}
\begin{longtable}{@{}p{3.5cm} p{3.2cm} p{3.2cm} p{3.0cm}@{}}
\toprule
\rowcolor{cvBlue}
\hdc{मानदंड} & \hdc{आरम्भिक (1)} & \hdc{अच्छा (2--3)} & \hdc{उत्कृष्ट (4)} \\
\midrule
\endhead
डेटा-संग्रह की गुणवत्ता & बहुत कम/अव्यवस्थित डेटा & पर्याप्त, कुछ त्रुटियाँ &
  पर्याप्त, स्वच्छ एवं सुव्यवस्थित \\
विश्लेषण एवं ग्राफ़ & ग्राफ़ अनुपस्थित/अस्पष्ट & सही ग्राफ़, सीमित व्याख्या &
  स्पष्ट ग्राफ़ एवं सही व्याख्या \\
पूर्वानुमान/निष्कर्ष & निष्कर्ष नहीं & निष्कर्ष है पर अधूरा &
  डेटा-समर्थित स्पष्ट निष्कर्ष \\
सीमाओं की समझ & सीमाएँ नहीं बताईं & कुछ सीमाएँ & सीमाएँ एवं सुधार सुझाए \\
प्रस्तुति एवं टीम-कार्य & अस्पष्ट प्रस्तुति & ठीक प्रस्तुति &
  स्पष्ट, सहभागी प्रस्तुति \\
\bottomrule
\end{longtable}

\vspace{6pt}
\noindent\textbf{सुझाव:} मूल्यांकन में \emph{प्रक्रिया} (कैसे सोचा, कैसे डेटा जुटाया)
को \emph{परिणाम} जितना ही महत्व दें -- यही वैज्ञानिक दृष्टिकोण का मूल है।
